\PassOptionsToPackage{no-math}{fontspec}
\documentclass[UTF8,a4paper,zihao=5,fontset=none]{ctexart}

\usepackage{fontspec}
\usepackage{xeCJK}
\setmainfont{Times New Roman}
\setsansfont{Times New Roman}
\setmonofont{Times New Roman}
\setCJKmainfont[AutoFakeBold=2.0,AutoFakeSlant=0.2]{SimSun}
\setCJKsansfont[AutoFakeBold=2.0,AutoFakeSlant=0.2]{SimSun}
\setCJKmonofont[AutoFakeBold=2.0]{SimSun}
\XeTeXgenerateactualtext=1

\usepackage[a4paper,top=1.8cm,bottom=1.7cm,left=1.85cm,right=1.85cm,footskip=0.7cm]{geometry}
\usepackage{amsmath,amssymb,booktabs,tabularx,array,enumitem,setspace,caption}
\usepackage[numbers,sort&compress]{natbib}
\usepackage[hidelinks]{hyperref}
\usepackage{multicol,titlesec,indentfirst,needspace}

\hypersetup{hypertexnames=false,pdftitle={AJAE：冻结STU的五帧稠密输入可行性研究},pdfsubject={激光雷达异常分割研究方案与可证伪实验协议}}
\captionsetup{font=small,labelfont=bf,labelsep=quad,justification=centering,singlelinecheck=true,hypcap=false}
\setstretch{1.06}
\setlength{\parindent}{2em}
\setlength{\parskip}{0pt}
\setlength{\tabcolsep}{3.5pt}
\setlength{\columnsep}{0.60cm}
\setlist[enumerate]{leftmargin=2.25em,itemsep=0.16em,topsep=0.22em,parsep=0pt}
\setlist[itemize]{leftmargin=2.10em,itemsep=0.13em,topsep=0.22em,parsep=0pt}
\renewcommand{\tabularxcolumn}[1]{m{#1}}
\newcolumntype{Y}{>{\raggedright\arraybackslash}X}
\newcolumntype{L}[1]{>{\raggedright\arraybackslash}m{#1}}
\titleformat{\section}{\bfseries\fontsize{12pt}{15pt}\selectfont}{\thesection}{0.55em}{}
\titlespacing*{\section}{0pt}{0.85em}{0.35em}
\titleformat{\subsection}{\bfseries\fontsize{10.5pt}{13pt}\selectfont}{\thesubsection}{0.50em}{}
\titlespacing*{\subsection}{0pt}{0.55em}{0.20em}
\pagestyle{plain}
\raggedbottom
\emergencystretch=1em

\begin{document}
\special{pdf:minorversion 7}
\fontsize{9.5pt}{12.2pt}\selectfont

\begin{center}
  {\bfseries\fontsize{16pt}{20pt}\selectfont
  AJAE：冻结STU的五帧稠密输入可行性研究}\\[0.15em]
  {\fontsize{9.5pt}{12pt}\selectfont 研究方案与可证伪实验协议（schema 33）}
\end{center}

\section*{摘要}

未知道路障碍物在单帧激光雷达中常只有少量且受遮挡的回波。AJAE首先检验一个比新模型设计更基础的问题：连续五帧对齐到当前扫描后，冻结的STU能否直接从更稠密的输入中获得更好的异常证据。窗口定义为当前帧及其前四帧，前四帧通过自车位姿变换到当前帧坐标。五帧作为一张稠密伪扫描进入原始冻结STU，但系统只读取当前帧原有点的MaxLogit异常分数。主比较因而是原始单帧STU与五帧稠密STU，不训练新的单帧网络，也不进行重叠窗口分数融合。

实验按信息增益依次进行。几何实验先测量回波数、唯一体素覆盖和新增体素，并用PLY检查配准。正常稳定性实验再比较两种输入的MaxLogit分布、正常语义准确率和平均交并比。少量合成异常实验为每段连续序列固定一个异常世界，以世界为配对单位比较点级平均精确率。只有冻结STU不能直接适应稠密输入时，才考虑在逐帧冻结STU特征之后训练小型异常头。当前完成的是schema 33路线和实现修订；三个真实可行性实验均未执行，因此本文不报告性能结论。

\noindent\textbf{关键词：} 激光雷达；异常分割；多帧稠密化；冻结模型；异常代理

\section*{Abstract}

Unknown road obstacles may produce only sparse and occluded returns in a single LiDAR scan. AJAE first tests a question that precedes the design of any new trainable model: whether a frozen STU model can obtain stronger anomaly evidence from five scans aligned to the current scan. Each causal window contains the current scan and its four predecessors. The five scans are treated as one dense pseudo-scan, while MaxLogit scores are read only for points originating from the current scan. The primary comparison is therefore the released single-scan STU input against the five-scan dense STU input.

The study proceeds through three pre-training experiments. Geometry statistics and PLY files first verify actual spatial densification. Normal-only data then test whether the dense input shifts MaxLogit or degrades semantic recognition. Finally, a small set of continuously rendered synthetic anomaly worlds tests paired changes in point-level average precision. A trainable anomaly head is considered only if direct dense input is rejected. No schema-33 feasibility experiment has yet been run, so this document makes no performance claim.

\noindent\textbf{Keywords:} LiDAR; anomaly segmentation; multi-scan densification; frozen model; anomaly proxy

\section{研究问题}

STU提供连续激光雷达扫描、逐点异常标注、冻结模型和评价程序\cite{nekrasov2025stu,stuofficial}。多扫描语义分割研究说明相邻观测可以缓解单帧稀疏\cite{li2023memoryseg,wu2024taseg}，但已知类别分割的改善不能直接推出未知障碍物检测改善。AJAE当前的研究问题是：

\begin{quote}
五帧激光雷达扫描对齐到当前帧形成的稠密伪扫描，能否改善冻结STU的异常证据，同时不造成不可接受的正常输出退化？
\end{quote}

这个问题包含两个必要部分。异常信号增强只回答稠密观测是否有用；正常稳定性回答原本在单帧点云上训练的STU能否承受输入分布变化。只有两部分同时成立，直接稠密STU才构成可继续验证的方法。

schema 32曾预先建立B0、B1、B2、B3训练矩阵，用对称参考系处理五帧并为全部点输出，再融合重叠窗口分数。该设计在核心可行性尚未验证前引入了额外单帧网络、扫描隔离消融、联合点网络和复杂训练选择。schema 33停止这条活动路线。旧实现和旧机械检查不能作为当前方法的性能证据。

\section{在线五帧观测}

当前时刻为$t$，窗口为

\begin{equation}
\mathcal W_t=\{X_{t-4},X_{t-3},X_{t-2},X_{t-1},X_t\}.
\end{equation}

令$T_{W\leftarrow i}$表示第$i$帧激光雷达坐标到世界坐标的刚体变换。过去帧通过

\begin{equation}
T_{t\leftarrow i}=T_{W\leftarrow t}^{-1}T_{W\leftarrow i}
\end{equation}

变换到当前扫描坐标，形成

\begin{equation}
\widetilde X_t=\bigcup_{i=t-4}^{t}T_{t\leftarrow i}X_i.
\end{equation}

该坐标系就是车辆当前传感器坐标，不是人为构造的中间参考空间。第$t$帧作为输出时刻具有明确作用，但五帧都作为STU的输入信息。系统只为$X_t$中的可见点返回异常分数。窗口从$t$移动到$t+1$时只输出新当前帧，因而无需平均一个点在多个窗口中的分数。

五帧位姿只补偿自车运动。动态车辆或行人在当前坐标下可能形成拖影。F1通过点云检查直接暴露这种现象，F2检验它是否使正常输出失稳；当前不为移动正常目标引入专用损失或模型分支。

\section{冻结STU比较}

\subsection{单帧输入}

单帧条件使用发布的STU预处理、冻结权重和MaxLogit异常分数：

\begin{equation}
X_t\xrightarrow{\operatorname{STU}_{\mathrm{frozen}}}S_t^{\mathrm{single}}.
\end{equation}

该条件是当前唯一基线。它不包含额外训练的单帧AJAE网络。

\subsection{稠密输入}

稠密条件将$\widetilde X_t$视为一张伪扫描，继续使用强度和扫描中心距离作为STU输入通道：

\begin{equation}
\widetilde X_t\xrightarrow{\operatorname{STU}_{\mathrm{frozen}}}S_{t-4:t}^{\mathrm{dense}}.
\end{equation}

随后按来源身份只选择$X_t$对应的行，得到$S_t^{\mathrm{dense}}$。单帧和稠密条件使用同一个模型、同一权重和同一分数定义，并在相同当前点上比较。五帧稠密输入会同时改变点密度、体素占用和扫描中心距离分布；F2必须把这些变化作为真实方法属性检验。

\section{数据划分}

\begin{table}[htbp]
\centering
\caption{schema 33的数据角色}
\begin{tabularx}{\textwidth}{L{2.2cm}L{2.6cm}Y}
\toprule
来源 & 范围 & 用途 \\
\midrule
正常开发 & \texttt{train/201}帧4--553 & F1几何、F2正常稳定性和F3少量合成异常；不更新参数 \\
正常确认 & \texttt{train/201}帧554--681 & 方法确定后一次性检查正常安全；此前不读取分数 \\
未来训练 & \texttt{train/206}帧0--448 & 仅在进入F4后生成完整虚拟序列并更新小型异常头 \\
真实异常 & 19条公开验证序列 & 方法和阈值固定后检验合成到真实的迁移 \\
隐藏测试 & 51条测试序列 & 公开验证支持方法后进行最终测试 \\
\bottomrule
\end{tabularx}
\end{table}

\texttt{train/201}帧0--3已确认存在内部精确重复，因此不进入开发范围。正常开发和正常确认是连续且不重叠的区间。确认区间不能参与输入选择、阈值选择或训练停止。

\section{递进实验}

\subsection{F1：几何稠密性}

F1遍历正常开发区间的546个合法窗口，不运行STU。每个窗口报告五帧与当前单帧的回波数之比，在0.05、0.10和0.20米体素下报告唯一体素数之比及当前帧没有覆盖的新体素数。当前帧8、198、387和553导出带扫描年龄的PLY，用于检查位姿方向、表面重合和动态目标拖影。

点数增加本身不能支持异常性能改善。若当前帧经过对齐后发生变化，或PLY显示普遍配准错误，必须停止后续实验。

\subsection{F2：正常稳定性}

F2在开发区间预先均匀选定的24个当前帧上运行单帧和稠密STU。两者只比较相同当前点。主要统计为MaxLogit中位数和第95百分位的变化、19类正常语义准确率和平均交并比。

稠密STU只有在MaxLogit中位数增加不超过0.02、第95百分位增加不超过0.05、正常语义准确率下降不超过0.02时视为稳定。平均交并比用于解释哪些类别变化，不新增事后门槛。

\subsection{F3：少量合成异常信号}

F3在正常开发区间选取8个互不重叠的28帧片段。每个片段采样一个只含异常代理的固定世界，随后连续渲染全部帧。从每段第5帧开始产生24个候选在线输出，同一异常在相邻窗口中保持同一世界位置、形状、材质和射线观测。只有在2.5--50.0米范围内至少含5个异常点的当前帧才进入点级平均精确率计算；结果同时记录候选帧数和实际纳入帧数。

两种条件在2.5--50.0米的相同当前点上计算点级平均精确率。每个世界产生一对结果，再以世界为配对单位执行10,000次自助重采样。只有F2通过，且

\begin{equation}
\Delta AP=AP_{\mathrm{dense}}-AP_{\mathrm{single}}
\end{equation}

的95\%区间下界大于0，直接稠密STU才得到合成开发证据支持。

\subsection{F4：最小训练分支}

F4只在F2或F3否定直接输入后启动。第一候选方法让五帧分别以原生单帧形式通过冻结STU，再把冻结特征和对齐后的几何送入一个小型异常头。训练样本来自多个在\texttt{train/206}上完整渲染的虚拟异常序列；相邻窗口复用同一物理帧。每轮可以从冻结窗口池抽样，避免重复计算全部窗口。

F1--F3完成前不规定F4的网络、损失和训练预算。大型Point Transformer只有在小型异常头被实际结果否定后才有讨论依据。

\section{评价、状态与限制}

合成异常的主指标为点级平均精确率，AUROC和FPR95为辅助指标。所有条件使用相同距离掩码、当前点身份和标签。移动正常目标仍按原始语义计为正常；只有主异常结果为正时才报告该子组。

当前权威协议为schema 33。F0的末帧坐标、成对冻结STU输入和F1--F3入口已经通过机械核验，当前状态是待执行的F1。当前没有schema 33真实几何统计、合成序列、训练检查点、正常确认结果、真实异常结果或隐藏测试结果。现阶段只能声称路线和实现已经收束，不能声称五帧输入改善了STU。

反事实渲染仍近似刚性扫描和首次回波，不覆盖扫描内运动畸变、束发散、多回波、多径及全部电子噪声。位姿对齐不能静止动态交通参与者。合成异常上的正向差异仍需公开真实异常序列确认，不能直接解释为真实道路泛化。

\section{预期产出}

若直接稠密STU得到支持，AJAE的核心产出将是一个无需训练的新输入构造及其正常安全和真实异常证据。若直接输入失败但小型异常头有效，AJAE的核心将转为冻结STU之后的多帧稠密特征适配。两条路线都必须由F1--F3的真实结果决定，而不能由预先设计的实验矩阵决定。

\noindent\textbf{数据可得性。} STU数据、评价代码和模型资源由数据集作者发布\cite{stuofficial}。本研究不重新发布原始扫描。

\noindent\textbf{伦理边界。} 本研究不新增道路采集，也不涉及人体或动物干预。若以后新增可识别图像或现场采集，应重新核对数据许可、隐私和机构审查要求。

\noindent\textbf{人工智能辅助。} 本文在代码审查、文字修订和排版中使用生成式人工智能辅助。研究者负责执行真实实验并审定科学结论；人工智能没有生成或替代实验数据。

\Needspace{12\baselineskip}
\phantomsection
\section*{参考文献}
\addcontentsline{toc}{section}{参考文献}
\begin{multicols}{2}
\fontsize{7.4pt}{8.5pt}\selectfont
\renewcommand{\bibsection}{}
\begin{thebibliography}{99}
\setlength{\itemsep}{0.25em}
\setlength{\parskip}{0pt}

\bibitem{nekrasov2025stu}
A. Nekrasov, M. Burdorf, S. Worrall, B. Leibe, and J. S. B. Perez. Spotting the Unexpected (STU): A 3D LiDAR Dataset for Anomaly Segmentation in Autonomous Driving. \textit{CVPR}, 2025, pp. 11875--11885.
\href{https://openaccess.thecvf.com/content/CVPR2025/html/Nekrasov_Spotting_the_Unexpected_STU_A_3D_LiDAR_Dataset_for_Anomaly_CVPR_2025_paper.html}{Official paper page}.

\bibitem{stuofficial}
A. Nekrasov et al. STU Dataset: Data, Evaluation Code, Checkpoints, and Test Submission Instructions. Official repository, 2025--2026.
\href{https://github.com/kumuji/stu_dataset}{Official repository}.

\bibitem{li2023memoryseg}
E. Li, S. Casas, and R. Urtasun. MemorySeg: Online LiDAR Semantic Segmentation with a Latent Memory. \textit{ICCV}, 2023, pp. 745--754.
\href{https://openaccess.thecvf.com/content/ICCV2023/html/Li_MemorySeg_Online_LiDAR_Semantic_Segmentation_ICCV_2023_paper.html}{Official paper page}.

\bibitem{wu2024taseg}
X. Wu et al. TASeg: Temporal Aggregation Network for LiDAR Semantic Segmentation. \textit{CVPR}, 2024, pp. 15311--15320.
\href{https://openaccess.thecvf.com/content/CVPR2024/html/Wu_TASeg_Temporal_Aggregation_Network_for_LiDAR_Semantic_Segmentation_CVPR_2024_paper.html}{Official paper page}.

\end{thebibliography}
\end{multicols}
\end{document}
